\documentclass[11pt]{article}
% ======================================================================
%  weekpreamble.tex — shared preamble for the weekly course summaries (EN)
%  Concise, readable, intuition-first. Same box style as the main doc.
%  Compiles with pdflatex.
% ======================================================================
\usepackage[utf8]{inputenc}
\usepackage[T1]{fontenc}
\usepackage[english]{babel}
\usepackage[a4paper,margin=2.2cm]{geometry}
\usepackage{amsmath,amssymb,amsthm,mathtools}
\usepackage{bm}
\usepackage{enumitem}
\usepackage{xcolor}
\usepackage[most]{tcolorbox}
\usepackage{microtype}

\emergencystretch=3em
\allowdisplaybreaks[1]
\setlength{\parindent}{0pt}
\setlength{\parskip}{0.45em}

% ---------- math macros (same names as the main document) ----------
\newcommand{\E}{\mathbb{E}}
\DeclareMathOperator{\Var}{Var}
\DeclareMathOperator{\Cov}{Cov}
\DeclareMathOperator{\Corr}{Corr}
\DeclareMathOperator*{\argmin}{arg\,min}
\DeclareMathOperator{\MSE}{MSE}
\DeclareMathOperator{\trace}{tr}
\newcommand{\Reals}{\mathbb{R}}
\newcommand{\Ints}{\mathbb{Z}}
\newcommand{\Comp}{\mathbb{C}}
\newcommand{\Nat}{\mathbb{N}}
\newcommand{\Prob}{\mathbb{P}}
\newcommand{\Tr}{^{\mathsf{T}}}
\newcommand{\Herm}{^{\mathsf{H}}}
\newcommand{\conj}[1]{\overline{#1}}
\newcommand{\abs}[1]{\left\lvert #1 \right\rvert}
\newcommand{\norm}[1]{\left\lVert #1 \right\rVert}
\newcommand{\ind}{\mathbf{1}}
\newcommand{\eps}{\varepsilon}
\newcommand{\diff}{\nabla}
\newcommand{\acvs}{\gamma}
\newcommand{\acf}{\rho}
\newcommand{\pacf}{\alpha}
\newcommand{\sdf}{\mathcal{S}}
\newcommand{\filt}{\mathcal{L}}
\newcommand{\Bsh}{B}
\newcommand{\ms}{\;\overset{\mathrm{ms}}{=}\;}
\newcommand{\toprob}{\xrightarrow{P}}
\newcommand{\todist}{\xrightarrow{d}}
\newcommand{\iid}{\overset{\text{iid}}{\sim}}

\setlist[itemize]{leftmargin=1.5em,topsep=2pt,itemsep=2pt}
\setlist[enumerate]{leftmargin=1.7em,topsep=2pt,itemsep=2pt}

% ---------- compact, titled (unnumbered) boxes ----------
\tcbset{wkbase/.style={enhanced,breakable,fonttitle=\bfseries,coltitle=white,
  boxrule=0.6pt,arc=1.2mm,left=2.2mm,right=2.2mm,top=1.1mm,bottom=1.1mm,
  before skip=6pt,after skip=6pt,
  attach boxed title to top left={yshift=-3.2mm,xshift=2.4mm},
  boxed title style={boxrule=0pt,arc=0.5mm}}}

\newtcolorbox{defn}[1]{wkbase, colback=blue!4!white, colframe=blue!58!black,
  colbacktitle=blue!58!black, title={Definition --- #1}, top=3mm}
\newtcolorbox{result}[1]{wkbase, colback=red!4!white, colframe=red!60!black,
  colbacktitle=red!60!black, title={#1}, top=3mm}
\newtcolorbox{intu}[1][Intuition]{wkbase, colback=yellow!12!white,
  colframe=orange!72!black, colbacktitle=orange!72!black, coltitle=white,
  title={#1}, top=3mm}
\newtcolorbox{retenir}[1][Key takeaways --- the week's reflexes]{wkbase,
  colback=green!5!white, colframe=green!48!black, colbacktitle=green!48!black,
  title={#1}, top=3mm}
\newtcolorbox{exmpl}[1][Worked example]{wkbase, colback=black!3!white,
  colframe=black!55!white, colbacktitle=black!55!white, title={#1}, top=3mm}

% ---------- week header ----------
\newcommand{\weekheader}[4]{%
  {\small\sffamily\bfseries\color{black!60} TIME SERIES~~$\cdot$~~MATH-342, EPFL}\par
  \vspace{2pt}
  {\LARGE\bfseries Week #1 --- #3\par}
  \vspace{1pt}
  {\itshape\color{black!55} #2}\par
  \vspace{6pt}
  \begin{tcolorbox}[enhanced,colback=blue!3!white,colframe=blue!45!black,
    arc=1.4mm,boxrule=0.7pt,left=3mm,right=3mm,top=2mm,bottom=2mm,before skip=2pt,after skip=10pt]
    \textbf{The big picture.}~ #4
  \end{tcolorbox}
}

\begin{document}
\weekheader{05}{19 March 2026}{Trends, Seasonality \& SARIMA}{Make a series stationary by removing trend and seasonality through differencing; ARIMA and SARMA encode this in the polynomials.}

Real series are almost never stationary: they \emph{drift} (a trend) and they \emph{repeat} (a season). The working model of the chapter is $X_t=\mu_t+Y_t$, with $\mu_t$ a deterministic time-dependent mean and $\{Y_t\}$ stationary. There are two ways to peel $\mu_t$ off: \textbf{estimate-and-subtract} (regress on $t$ and on seasonal dummies) or \textbf{difference it away} with $\diff=I-B$. Layering AR/MA structure onto both ordinary and seasonal lags gives the SARMA family; folding the differencing in front gives (S)ARIMA --- the Box--Jenkins workhorse.

\section*{Key definitions}

\begin{defn}{Mean-trend (signal-plus-noise) model}
$X_t=\mu_t+Y_t$, where $\mu_t$ is a deterministic, time-dependent mean (the \textbf{trend}) and $\{Y_t\}$ is zero-mean second-order stationary. Canonical cases: $\mu_t=a+bt$ (linear), $\mu_t$ a degree-$(q-1)$ polynomial, or $\mu_t=s_t$ \textbf{periodic} of period $s$ ($s_{t+s}=s_t$). The two can coexist: $X_t=a+bt+s_t+Y_t$.
\end{defn}

\begin{defn}{Difference operators $\diff$ and $\diff_{(s)}$}
Ordinary difference: $\diff=I-B$, $\ \diff X_t=X_t-X_{t-1}$ (with $BX_t=X_{t-1}$). It is linear, and
\[
\diff^{d}=(I-B)^{d}=\sum_{k=0}^{d}\binom{d}{k}(-1)^{k}B^{k},\qquad
\text{e.g. } \diff^{2}X_t=X_t-2X_{t-1}+X_{t-2}.
\]
\textbf{Seasonal} difference of period $s$: $\diff_{(s)}=I-B^{s}$, $\ \diff_{(s)}X_t=X_t-X_{t-s}$ (just two terms). Ordinary is the case $s=1$.
\end{defn}

\begin{defn}{Seasonal MA and AR (SMA, SAR)}
Only lags that are \emph{multiples of $s$} appear:
\[
\begin{aligned}
\mathrm{SMA}(Q)_s:&\quad X_t=\eps_t-\sum_{j=1}^{Q}\theta^{(s)}_j\eps_{t-js}=\Theta^{(s)}(B)\eps_t,\\
\mathrm{SAR}(P)_s:&\quad \Phi^{(s)}(B)X_t=\eps_t,\quad X_t=\eps_t+\sum_{j=1}^{P}\phi^{(s)}_jX_{t-js}.
\end{aligned}
\]
Polynomials: $\Theta^{(s)}(z)=1-\sum_j\theta^{(s)}_jz^{sj}$ and $\Phi^{(s)}(z)=1-\sum_j\phi^{(s)}_jz^{sj}$. Basic seasonal MA(1): $X_t=\eps_t-\theta\eps_{t-s}$.
\end{defn}

\begin{defn}{SARMA$(p,q)\times(P,Q)_s$}
Multiply an ordinary ARMA by a seasonal ARMA at period $s$:
\[
\Phi(B)\,\Phi^{(s)}(B)\,X_t=\Theta(B)\,\Theta^{(s)}(B)\,\eps_t .
\]
$\mathrm{SMA}(Q)_s=\mathrm{SARMA}(0,0)\times(0,Q)_s$ and $\mathrm{SAR}(P)_s=\mathrm{SARMA}(0,0)\times(P,0)_s$. Example, $\mathrm{SARMA}(0,1)\times(0,1)_{12}$: $\ X_t=(1-\theta B)\bigl(1-\theta^{(12)}B^{12}\bigr)\eps_t$.
\end{defn}

\begin{defn}{ARIMA$(p,d,q)$ and SARIMA$(p,d,q)\times(P,D,Q)_s$}
$\{X_t\}$ is \textbf{ARIMA$(p,d,q)$} (``integrated ARMA'') if its $d$-th difference is a causal, invertible ARMA$(p,q)$:
\[
\Phi(B)\,(1-B)^{d}X_t=\Theta(B)\,\eps_t ,\qquad\text{i.e. } Z_t=\diff^{d}X_t \text{ is ARMA}(p,q).
\]
The ``I'' is \emph{integrated}: recover $X_t$ from $Z_t$ by summing $d$ times. The full \textbf{SARIMA} folds in both an ordinary difference $d$ and a seasonal difference $D$:
\[
\Phi(B)\,\Phi^{(s)}(B)\,(1-B)^{d}(1-B^{s})^{D}\,X_t=\Theta(B)\,\Theta^{(s)}(B)\,\eps_t .
\]
\end{defn}

\begin{defn}{Random walk, IMA$(1,1)$, ARI$(1,1)$}
\textbf{Random walk} $X_t=X_{t-1}+\eps_t$ is the AR(1) at $\phi=1$: \emph{not} stationary, but $\diff X_t=\eps_t$ is white noise, so it is ARIMA$(0,1,0)$. Two minimal one-step-integrated models, both with non-stationary $\{Y_t\}$:
\begin{itemize}
\item \textbf{IMA$(1,1)$}=ARIMA$(0,1,1)$: $\ Y_t=Y_{t-1}+\eps_t-\theta\eps_{t-1}$ (i.e. $\diff Y_t=\eps_t-\theta\eps_{t-1}$);
\item \textbf{ARI$(1,1)$}=ARIMA$(1,1,0)$: $\ Y_t-Y_{t-1}=\phi(Y_{t-1}-Y_{t-2})+\eps_t$, $\abs\phi<1$.
\end{itemize}
\end{defn}

\begin{intu}[Subtract vs.\ difference]
\textbf{Regress-and-subtract} ($a+bt$, or $\sin/\cos$ harmonics) gives an interpretable deterministic trend, but you must \emph{know its form}. \textbf{Differencing} is model-free and robust: $\diff$ kills any linear drift, $\diff^{q}$ any degree-$(q-1)$ polynomial, $\diff_{(s)}$ any period-$s$ season --- at the price of slightly distorting the noise ($\diff Y_t$ replaces $Y_t$) and risking \emph{over}-differencing.
\end{intu}

\section*{Key results \& formulas}

\begin{result}{Differencing removes a polynomial trend}
If $X_t=a+bt+Y_t$ with $\{Y_t\}$ stationary, then $\diff X_t=b+\diff Y_t$: the trend collapses to the constant $b$, and $\diff Y_t$ is again stationary. More generally, if $\mu_t$ is a degree-$(q-1)$ polynomial then $\diff^{q}\mu_t\equiv0$, so
\[
\diff^{q}X_t=\diff^{q}Y_t=\sum_{k=0}^{q}\binom{q}{k}(-1)^{k}Y_{t-k}.
\]
\textit{Idea:} each $\diff$ lowers a polynomial's degree by one; linearity passes through to $Y$.
\end{result}

\begin{intu}[How many differences?]
$\diff$ lowers a polynomial degree by exactly one. So \textbf{linear drift $\Rightarrow d=1$}, \textbf{quadratic $\Rightarrow d=2$}: choose $d=q$ to flatten a degree-$(q-1)$ trend. A seasonal pattern of period $s$ needs \emph{one} seasonal difference $\diff_{(s)}$ ($D=1$).
\end{intu}

\begin{result}{Theorem --- $\diff\diff_{(12)}$ flattens trend + season}
Let $X_t=a+bt+s_t+Y_t$ with $s_{t+12}=s_t$ and $\{Y_t\}$ zero-mean stationary. Apply the seasonal difference first ($\diff_{(12)}X_t=12b+\diff_{(12)}Y_t$ kills $s_t$ and turns $bt$ into a constant), then $\diff$ kills the constant:
\[
\diff\diff_{(12)}X_t=\diff\diff_{(12)}Y_t=Y_t-Y_{t-1}-Y_{t-12}+Y_{t-13}.
\]
This is a fixed linear combination of stationary $Y$'s: mean $0$, covariance a function of $\tau$ only $\Rightarrow$ weakly stationary. (This is the standard ``$\diff_{12}$ then $\diff$'' pre-processing for monthly data, e.g.\ the Mauna Loa CO$_2$ series.)
\end{result}

\begin{intu}[Multiplicative season needs a different recipe]
If trend and season \emph{multiply}, $X_t=(a+bt)s_t+Y_t$ (amplitude grows with level), then $\diff\diff_{(12)}$ is \emph{wrong}: here a single \emph{double seasonal} difference works, $\diff_{(12)}^{2}X_t=Y_t-2Y_{t-12}+Y_{t-24}$, again stationary.
\end{intu}

\begin{result}{ACVS of the seasonal MA(1)}
For $X_t=\eps_t-\theta\eps_{t-s}$ (white-noise variance $\sigma_\eps^2$), only three lags survive:
\[
\acvs_0=(1+\theta^{2})\sigma_\eps^2,\quad \acvs_{\pm s}=-\theta\,\sigma_\eps^2,\quad \acvs_\tau=0\ \text{else}
\qquad\Rightarrow\qquad \acf_{\pm s}=\frac{-\theta}{1+\theta^{2}}.
\]
\textit{Idea:} white-noise filter rule $\acvs_\tau=\sigma_\eps^2\sum_j a_ja_{j+\abs\tau}$ on taps $a_0=1,a_s=-\theta$. A single ACF spike at lag $s$ is the fingerprint of an $\mathrm{SMA}(1)_s$.
\end{result}

\begin{result}{Random walk and integrated growth (the difference-me signal)}
For the random walk $X_t=\sum_{j=1}^{t}\eps_j$ ($X_0=0$): $\E[X_t]=0$, $\Var(X_t)=t\sigma_\eps^2$, $\Cov(X_t,X_{t+\tau})=t\sigma_\eps^2$ --- all depend on $t$, so $\{X_t\}$ is \emph{not} stationary, yet $\diff X_t=\eps_t$ is white noise. Same diagnostic for IMA$(1,1)$ ($\theta\neq1$):
\[
\Var(Y_t)=\sigma_\eps^2\bigl[1+(1-\theta)^{2}(t-1)+\theta^{2}\bigr],\qquad \Corr(Y_t,Y_{t+\tau})\to1.
\]
ARI$(1,1)$ behaves the same way. \textbf{Variance growing linearly in $t$ and nearby values almost perfectly correlated} is the signature that a difference is needed.
\end{result}

\begin{result}{Theorem (pitfall) --- $\diff_{(s)}\neq\diff^{s}$, and over-differencing}
$\diff_{(s)}=I-B^{s}$ subtracts the value $s$ steps back (two terms, removes a period-$s$ \emph{season}); $\diff^{s}=(I-B)^{s}$ is an $(s+1)$-term binomial sum (removes a degree-$(s-1)$ \emph{polynomial}). They are not the same. Differencing an already-stationary series injects a $(1-B)$ factor into the MA part: a non-invertible MA(1) with $\acf_1\to-\tfrac12$. Seeing $\acf_1\approx-\tfrac12$ after differencing means \emph{difference one time fewer}.
\end{result}

\begin{exmpl}[Seasonal ARMA $(1-0.7B^{2})X_t=(1-0.3B^{2})\eps_t$]
Write $\phi=0.7,\theta=0.3$. AR roots have modulus $>1$ $\Rightarrow$ causal; expand the geometric series in $B^{2}$. Only \emph{even} lags survive: $a_0=1$, $a_j=0$ for odd $j$, and $a_j=\phi^{(j-2)/2}(\phi-\theta)$ for even $j\ge2$. Hence odd-lag correlations vanish and the ACF decays geometrically at rate $\phi$ along even lags: $\acf_2\approx0.47$, $\acf_4\approx0.33$, $\acf_6\approx0.23$ --- the correlogram shows spikes only at $\tau=2,4,6,\dots$ decaying like $0.7^{\tau/2}$.
\end{exmpl}

\section*{Intuition}

\begin{intu}[The whole chapter in one line]
Strip non-stationarity \emph{before} modelling: $\diff^{d}$ for trend, $\diff_{(s)}^{D}$ for season; then fit a SARMA on the (now stationary) result. SARIMA bundles all of this into four polynomials plus two differencing exponents.
\end{intu}

\begin{intu}[Reading a correlogram]
\textbf{Ordinary} ARMA structure lives at small lags $1,2,3,\dots$; \textbf{seasonal} structure clusters around $s,2s,\dots$ (with small side-lobes at $s\pm1$ from cross terms). A slowly decaying sample ACF $\Rightarrow$ a unit root $\Rightarrow$ difference. A lone $-\tfrac12$ at lag 1 $\Rightarrow$ you over-differenced.
\end{intu}

\section*{Key takeaways}

\begin{retenir}
\begin{itemize}
\item \textbf{Recipe.} $X_t=\mu_t+Y_t$: remove $\mu_t$ by regression (need the form) \emph{or} differencing (model-free). For monthly data with linear drift: $\diff\diff_{(12)}$ (seasonal first, then ordinary).
\item \textbf{Pick $d,D$.} $d=q$ flattens a degree-$(q-1)$ polynomial trend; $D=1$ removes a period-$s$ season (rarely $D=2$). Stop when the mean is stable, variance constant, and the ACF decays fast.
\item \textbf{Know the polynomials.} ARIMA: $\Phi(B)(1-B)^{d}X_t=\Theta(B)\eps_t$. SARIMA multiplies in seasonal factors $\Phi^{(s)},\Theta^{(s)}$ and $(1-B^{s})^{D}$. Expand by treating $B$ as a formal variable; seasonal factors contribute only multiples of $s$, cross terms give lags $s\pm1$.
\item \textbf{Diagnostics of integration.} Variance growing $\propto t$ and near-unit correlation of nearby points $\Rightarrow$ difference. Random walk, IMA$(1,1)$, ARI$(1,1)$ all show this.
\item \textbf{Traps.} $\diff_{(s)}\neq\diff^{s}$ (season vs.\ polynomial). Over-differencing creates a non-invertible MA unit root with $\acf_1\to-\tfrac12$.
\end{itemize}
\end{retenir}

\end{document}
